\section*{Bab \ref{ch:g6:earth-a-planet} --- Bumi, Sebuah Planet dalam Tata Surya}

\begin{solution}{exo:g6:earth-a-planet:1}
\underline{Merkurius}, \underline{Venus}, \underline{Bumi},
\underline{Mars}, Jupiter, Saturnus, Uranus, Neptunus. Planet-planet \omterm{def:g3:stones-pebbles-sand:rock}{batuan}
adalah keempat planet yang pertama.
\end{solution}

\begin{solution}{exo:g6:earth-a-planet:2}
``\omterm{def:g6:earth-a-planet:planet-types}{Planet batuan} itu kecil dan padat, dengan permukaan yang padat. \omterm{def:g6:earth-a-planet:planet-types}{Planet
raksasa} itu sangat besar, dengan massa jenis yang rendah, sebagian besar
tersusun dari gas dan cairan.''
\end{solution}

\begin{solution}{exo:g6:earth-a-planet:3}
Venus (12\,104~km, dibandingkan dengan 12\,756~km untuk Bumi). Planet yang
paling besar adalah Jupiter (142\,984~km); yang paling kecil adalah
Merkurius (4\,879~km).
\end{solution}

\begin{solution}{exo:g6:earth-a-planet:4}
(a) Salah: di Venus tidak ada air cair; planet itu terlalu panas. (b) Benar
(rata-rata $-65$~$^\circ$C, dibandingkan dengan 15~$^\circ$C). (c) Salah:
keempat planet yang paling dekat dengan Matahari adalah \omterm{def:g6:earth-a-planet:planet-types}{planet batuan};
planet-planet raksasa adalah keempat planet yang paling jauh. (d) Benar:
kira-kira 0.7~gram per sentimeter kubik.
\end{solution}

\begin{solution}{exo:g6:earth-a-planet:5}
$4879 \div 12756 \approx 0.382$, yaitu 38~\%: lebar Merkurius sedikit lebih
dari sepertiga lebar Bumi.
\end{solution}

\begin{solution}{exo:g6:earth-a-planet:6}
$120536 \div 12756 \approx 9.4$: kira-kira 9 Bumi. Saturnus adalah \omterm{def:g6:earth-a-planet:planet-types}{planet raksasa}.
\end{solution}

\begin{solution}{exo:g6:earth-a-planet:7}
Planet~A adalah \omterm{def:g6:earth-a-planet:planet-types}{planet raksasa}: planet itu sangat besar dan massa jenisnya,
1.1~gram per sentimeter kubik, rendah, seperti massa jenis planet-planet
raksasa di Tata Surya (antara 0.7 dan 1.7). Planet~B adalah \omterm{def:g6:earth-a-planet:planet-types}{planet batuan}:
planet itu kecil dan massa jenisnya, 5.0, tinggi, seperti massa jenis
planet-planet \omterm{def:g3:stones-pebbles-sand:rock}{batuan} (antara 3.9 dan 5.5).
\end{solution}

\begin{solution}{exo:g6:earth-a-planet:8}
(a) $464 - 15 = 449$~derajat lebih panas. (b) $15 - (-65) = 80$~derajat lebih dingin.
\end{solution}

\begin{solution}{exo:g6:earth-a-planet:9}
Venus mempunyai atmosfer karbon dioksida yang sangat tebal, yang menahan
panas Matahari seperti selimut tebal dan tidak membiarkannya lepas.
Merkurius, yang lebih dekat dengan Matahari, hampir tidak berudara, jadi
planet itu tidak dapat menyimpan panasnya.
\end{solution}

\begin{solution}{exo:g6:earth-a-planet:10}
$228 - 150 = 78$~juta kilometer lebih jauh. $78 \div 150 = 0.52$: jarak
Mars 52~\% lebih besar daripada jarak Bumi.
\end{solution}

\begin{solution}{exo:g6:earth-a-planet:11}
Miliaran tahun yang lalu, Mars mempunyai \omterm{def:g1:landscapes-around-us:water}{sungai}, \omterm{def:g1:landscapes-around-us:water}{danau}, dan \omterm{def:g5:water-shapes-the-land:flood}{banjir}: planet
itu lebih hangat dan lebih basah. Sekarang airnya membeku, di dalam es
tudung \omterm{def:g3:continents-and-oceans:poles}{kutub} dan di bawah tanah dekat \omterm{def:g3:continents-and-oceans:poles}{kutub}. Air cair tidak dapat bertahan
di permukaannya sekarang karena Mars sangat dingin (rata-rata
$-65$~$^\circ$C) dan karena udaranya begitu tipis sehingga air cepat
membeku atau berubah menjadi uap.
\end{solution}

\begin{solution}{exo:g6:earth-a-planet:12}
(a) $8.9 \div 9.8 \approx 0.908$, yaitu 91~\%. (b) Seperti kembaran:
ukurannya hampir sama (95~\% diameter Bumi), tarikan gravitasinya hampir
sama, dan Venus juga \omterm{def:g6:earth-a-planet:planet-types}{planet batuan}. Tidak seperti kembaran: Venus jauh lebih
panas (464~$^\circ$C dibandingkan dengan 15~$^\circ$C), udaranya jauh lebih
tebal dan terdiri atas karbon dioksida dengan \omterm{def:g4:clouds-rain-snow:cloud}{awan} asam, dan di sana tidak
ada air cair.
\end{solution}

\begin{solution}{pb:g6:earth-a-planet:1}
\textbf{1.} Venus (108~juta~km), Bumi (150), Mars (228).

\textbf{2.} $150 - 108 = 42$: Venus 42~juta kilometer lebih dekat.
$228 - 150 = 78$: Mars 78~juta kilometer lebih jauh.

\textbf{3.} $6792 \div 12756 \approx 0.532$, yaitu 53~\%: lebar Mars
kira-kira separuh lebar Bumi.

\textbf{4.} Ketiganya \omterm{def:g6:earth-a-planet:planet-types}{planet batuan}: massa jenisnya (5.2, 5.5, dan 3.9
gram per sentimeter kubik) tinggi, dan diameternya kecil dibandingkan
dengan diameter planet-planet raksasa.

\textbf{5.} Udara yang paling tebal ada di Venus (92~bar); yang paling
tipis di Mars (0.01~bar).

\textbf{6.} $1 \div 0.01 = 100$: udara Mars menekan seratus kali lebih
lemah daripada udara Bumi.

\textbf{7.} Udaranya yang tebal, yang terdiri atas karbon dioksida, menahan
panas Matahari dan tidak membiarkannya lepas, seperti selimut tebal.

\textbf{8.} $92 \div 1 = 92$: di tanah Venus, udara menekan
\textbf{92 kali} lebih kuat daripada di tanah Bumi. Wahana harus mempunyai
dinding yang sangat kuat agar tidak remuk (dan juga harus tahan terhadap
panas dan \omterm{def:g4:clouds-rain-snow:cloud}{awan} asam).

\textbf{9.} Suhu rata-rata Bumi, 15~$^\circ$C, berada di antara
0~$^\circ$C dan 100~$^\circ$C: air tidak membeku dan tidak mendidih, jadi
air tetap cair.

\textbf{10.} Di Venus, pada 464~$^\circ$C, jauh di atas 100~$^\circ$C,
air itu akan langsung mendidih habis dan berubah menjadi uap. Di Mars, pada
$-65$~$^\circ$C, jauh di bawah 0~$^\circ$C, air itu akan membeku menjadi es
(dan perlahan-lahan berubah menjadi uap, karena udaranya begitu tipis).

\textbf{11.} Dahulu kala, Mars lebih hangat dan lebih basah daripada
sekarang: air cair dapat mengalir di permukaannya. \omterm{def:g5:weather-and-climate:climate}{Iklimnya} telah sangat
berubah: sekarang Mars adalah \omterm{def:g5:weather-and-climate:desert}{gurun} yang dingin dan kering.

\textbf{12.} Misalnya: ``Wahana-wahana yang baik, di sini, di Bumi, kami
mempunyai air cair, di \omterm{def:g2:water-on-earth:ocean}{samudra}, \omterm{def:g1:landscapes-around-us:water}{sungai}, dan \omterm{def:g1:landscapes-around-us:water}{danau}. Suhu kami sedang,
rata-rata 15~$^\circ$C, di antara titik beku dan titik didih air. Kami
berada pada jarak yang tepat dari Matahari, tidak terlalu dekat dan tidak
terlalu jauh. Atmosfer kami menjaga kami tetap hangat, memberi makhluk hidup
gas yang mereka perlukan, dan membuat air tetap cair. Itulah sebabnya
kehidupan mungkin ada di sini.''
\end{solution}
